\section{Experiments}
This section provides the experimental procedures for each layer of the stack. With the goal of reproducibility and comparison, all experiments used a SPSA fixed random seed ($\texttt{seed}=42$), ensuring no variance between runs. 
Considering the variational principle, which guarantees $E_{VQE} \geq E_{FCI}$, results are reported as the "best physical energy", the lowest energy observed during 
the optimization remaining at or above the FCI. This addresses the HWE ansatz's limitation, as it does not conserve particle number, leading SPSA to push $\theta$
into exploring unphysical Hilbert space sectors (as discussed in Section~\ref{sec:hwe_limitation}, \parencite{kandala2017hardware}).

\subsection{Experiment 1: Distributed Statevector VQE (Simulator)}
This is the primary validation experiment of this stack, with the goal being to validate the distributed statevector evaluation path to return 
physically correct ground-state energy estimates and then to compare the distributed path's accuracy against the serial baseline's results.

In this experiment, each benchmark molecule from Table \ref{table:molecules} was configured with following identical variables:

\begin{itemize}
    \item \textbf{Backend}: Statevector simulator (exact, noiseless)  
    \item \textbf{Ansatz}: HWE-adaptive ansatz (EfficientSU2, full entanglement, auto selected reps). 
    \item \textbf{MPI ranks}: $P = 2$, the minimum rank count required to enable Pauli term distribution.
    \item \textbf{Optimizer}: SPSA with dimension scaled step size  ($a= 0.628 / \sqrt{n_{\text{params}} / 8}$). 
    \item \textbf{Seed}:42, remains fixed across all experiments to allow for reproducibility.
    \item \textbf{Max iterations}: $\max(200,\; n_{\text{params}} \times 8)$, scaling with problem dimension. 
    \item \textbf{Convergence}: Energy spread over the last 10 iterations $< 1.6 \times 10^{-3}$  Ha.         
    \item \textbf{Initialization}: Parameters drawn from $\mathcal{U}(-0.1, 0.1)$, keeps the initial state near the Hartree-Fock reference $|0\ldots0\rangle$. 
\end{itemize}   

% During iterations of the MPI workflow, subsets of the Pauli terms of the molecular Hamiltonian were distributed across MPI ranks, where each rank independently 
% builds the full $2^n$-dimensional statevector, which makes the stack's main bottleneck, costing $O(2^n)$ each iteration. Then in parallel, each rank uses their statevector to evaluates its local subset of $\langle\psi(\theta)|O_i|\psi(\theta)\rangle$ Pauli terms. 
% After evaluation is completed, a global \texttt{MPI\_Allreduce(SUM)} combines the partial energies from each rank. Lastly, the masking metric $M = T_{\text{accel}} / T_{\text{comm}}$  
% gets recorded one time per iteration.

A control experiment using an identical SPSA parameters with slightly differing in HWE ansatz reps calculation and convergence criteria within execution paths, leading to a difference of 72 vs. 96 parameters for $LiH$, making the serial baseline marginally faster than LiH (therefore, the reported speedup is conservative). 
This VQE was executed within a single Python process without using MPI ($P=1$), done to provide the $T_{\text{serial}}$ value. 
Note that the serial baseline may converge at a different iteration count than the distributed stack, as the convergence criterion can trigger earlier/later depending on the SPSA trajectory in each execution path, caused by the floating point errors that occur 
in \texttt{MPI\_Allreduce} reduction. 

\subsection{Experiment 2: Scaling Analyses}
The objective of this experiment is primarily to measure both Strong Scaling and Weak Scaling analyses across  $P \in \{1,2,4,8\}$ MPI ranks. 
Both scaling analysis are repeated with GPU acceleration within Experiment 4.

All four benchmark molecules (Table~\ref{table:molecules}) were run at full iteration counts for each rank 
count $P \in \{1,2,4,8\}$ (as powers of 2 is standard for scaling studies, providing straightforward mapping to binary tree reduction), all within one Docker container on a single host. 
Each configuration was executed once, as the deterministic statevector simulator with the fixed 
random seed ($\texttt{seed}=42$) eliminates any variance between runs. $H_{2}O$ (14 qubits, 1086 Pauli terms, 896 iterations) was selected as the main 
efficiency benchmark due to having the largest number of distributable Pauli terms. For each run, the total 
wall-clock time, the per-iteration time, and MPI communication time get recorded and parallel efficiency then computed as:
\begin{equation}
    E(P) = \frac{T_1}{P \times T_P}
\end{equation}
where $T_1$ is the determined wall-clock time at $P= 1$.

\subsubsection{Weak Scaling}
At each rank count, a molecule is selected so that the per-rank workload (the Pauli terms/rank) scales together, as stated within Section~\ref{sec:methodology}.   
Each molecular configuration is run 10 iterations, with all other configuration parameters remaining identical to those used 
in Experiment 1. This scaling strategy uses 10 iterations instead of full runs as this provides sufficient runs to measure the per-iteration timing. 
In this experiment, the stack's total wall-clock time, per-iteration time, per-iteration MPI communication time, and the median masking metric $M$ were recorded. 

\subsection{Experiment 3: IBM Quantum QPU Validation}   
This experiment is the real hardware validation, demonstrating the end-to-end middleware functionality when using a real quantum 
processor (QPU) instead of a simulator. This experiment validates the stack's Python EstimatorV2 integration, circuit transpilation pipeline, job submission, retrieving results, SPSA parameter updates,  
and the MPI coordination with cloud based QPU results. This experiment's goal is not accuracy, as 10 iterations is not sufficient to converge to a ground state energy. 

The molecule $H_{2}$ was selected for the IBM QPU experiment due to it requiring the least amount of qubits,
facing least complications with qubit decoherency. If the stack can function as expected for $H_2$, the 
middleware is validated and indicates accuracy will improve with more QPU iterations and error mitigation (which will be considered for future work).
This experiment uses an identical HWE-adaptive ansatz, SPSA hyperparameters, 
and near zero initialization as Experiment 1. However, this specific VQE was configured with:
\begin{itemize}
    \item \textbf{Backend}: IBM \texttt{ibm\_marrakesh} (using 4 of 156 physical qubits, Heron processor)\parencite{rallis2025interfacingquantumcomputingsystems}
    \item \textbf{MPI ranks}: $P=2$ (rank 0 manages QPU submission)
    \item \textbf{Shots}: 4096 per circuit (chosen to balance statistical precision of $1/\sqrt{4096} \approx 0.016\;E_h$ against the queue time on the IBM open access plan)
    \item \textbf{Error mitigation}: T-REx (resilience level 1, default EstimatorV2 settings)\parencite{cai2023quantumerrormitigation}
    \item \textbf{Max iterations}: 10 (constrained by the queue limits of IBM open access plan)
    \item \textbf{PUB batching}: Both $E(\theta^+)$ and $E(\theta^-)$ were submitted to SPSA as 2 Primitive Unified Blocs (PUBs) within a single EstimatorV2 job, reducing the QPU RTT wait by 50\%. 
\end{itemize}

The ansatz circuit uses logical/abstract qubits. Here, the ansatz was transpiled once against the physical QPU backend's specific coupling map by using Qiskit's preset pass 
manager (optimization level=1), then observable was then remapped to the physical qubit layout and then the one-time mapped circuit and observable layout are locally cached, 
as the circuit structure will not change and rebuilding will waste precious time. On IBM's 
open access plan, Sessions (allowing multiple jobs to share a QPU reservation) are not supported, so each iteration submits an independent job through the 
sessionless EstimatorV2 mode and then must re-enter the queue.    

The QPU experiment was executed on March 26, 2026 against the QPU backend calibration data that was valid at the time of the remote
job submission. Results for the IBM experiment are presented in Tables~\ref{table:ibm_results} and~\ref{table:ibm_iterations}.

\subsection{Experiment 4: GPU-Accelerated Statevector}
This experiment evaluates the effect of GPU acceleration on the distributed VQE pipeline, comparing the GPU 
backend against the CPU-only results from Experiments 1 and 2. This will determine if the GPU acceleration will bypass the statevector build bottleneck as predicted by \textcite{bayraktar2023cuquantum}. 

Within two different Docker configurations, all four benchmark molecules were run at $P \in \{1,2,4,8\}$ MPI ranks with GPU acceleration enabled, using the same HWE ansatz, SPSA hyperparameters, seed=42, and convergence criteria as Experiment 1 (all Python code remains unchanged). 
The GPU statevector configuration uses Qiskit Aer's \texttt{AerSimulator(method='statevector', device='GPU')}, which internally assigns to NVIDIA's cuStateVec library for GPU-resident statevector 
construction and gate operations. The Docker build configurations were tested for an evaluation of the effect of the cuStateVec integration on the stack's speed:
\begin{itemize}
    \item \textbf{GPU Aer (CUDA thrust)}: Qiskit Aer built with the CUDA thrust backend for basic GPU parallelism, with GPU-accelerated gate operations and the statevector management uses the default, unoptimized Aer CUDA kernels.
    \item \textbf{GPU Aer (cuStateVec)}: Qiskit Aer built at compile time with cuStateVec linkage from the NVIDIA cuQuantum SDK~\parencite{bayraktar2023cuquantum}, keeping the full statevector in GPU memory with all gate operations and measurement handled by cuStateVec's optimized kernels.
\end{itemize}

Both configurations use the same Python code path; the difference is found within the compiled Qiskit Aer backend linked at Docker build time.

The GPU hardware used was an NVIDIA GeForce GTX 1650 Mobile running CUDA 12.6, representing an accessible, lower bound on GPU-accelerated performance as described in Section ~\ref{sec:methodology}.   
% as the superior A100 HPC grade hardware
% would provide the stack with a substantially greater speedup. However, this experiment will still prove the proposed architectural hypothesis of this paper, 
% that GPU acceleration reduces time spent on the local serial statevector build enough to 
% continue scaling the MPI distribution. Future work with an A100 GPU that has 16x more bandwidth will result substantial improvement.  

For each configuration, total wall-clock time, per-iteration time, and speedup relative to the serial baseline 
were recorded. Strong scaling and weak scaling analyses were repeated with GPU acceleration to compare against 
the CPU-only scaling trials from Experiment 2.

\subsection{Experiment 5: Checkpoint Resilience}
This experiment is conducted to validate the checkpoint restart mechanism of the stack, which verifies that 
the optimization can resume from a previously saved $\theta$ state without a loss of convergence progress. Checkpoints 
are saved every 5 iterations (hardcoded in the file \texttt{interface.py}), along with a rolling retention of the 5 most recent 
checkpoint files logged per molecule. This interval balances recovery granularity (at most 5 iterations 
of work lost on failure) against I/O overhead when tested with less complex molecules within our benchmark. 

The molecule $H_{2}$ was selected in this experiment due to its low complexity and fast per iteration time, making it a 
practical choice for repeated restart testing in this stack. The 7-layer diagnostic suite (\texttt{test\_layers\_run.py}) automatically executes this experiment:

\begin{enumerate}
    \item Run $H_{2}$ VQE for 5 iterations, then save a checkpoint of the current $\theta$ vector to an .npy file.
    \item Verify the checkpoint file exists and contains the correct parameter vector shape (24 for $H_{2}$).
    \item Restart VQE from the checkpoint, running 5 additional iterations (6--10).
    \item Verify the optimizer resumes from the correct SPSA schedule position ($k=6$) and that the restarted energy continues from the checkpointed trajectory.
\end{enumerate}

The experiment passes if several criteria are satisfied: if the checkpoint file exists with the correct parameter vector shape, the optimizer has
resumed from the SPSA schedule position $k=6$, as otherwise a reset to the initial values occur resulting in the optimizer to take large steps (destroying convergence), 
and lastly, the restart energy at iteration 6 is within 0.5\,$E_h$ of the energy saved at iteration 5. 
The tolerance of 0.5\,$E_h$ is to account for the random nature of SPSA perturbations while rejecting absolute failures such as a corrupted parameter vector of loaded zeros.

\subsection{Experiment 6: Stress Testing}
Two additional stress tests were conducted to validate the stack's robustness to less ideal conditions as the ones explored in Experiment 5. 
Both tests also are executed automatically as part of the 7-layer diagnostic suite (\texttt{make trial}).

\subsubsection{QPU Latency Spiking}
This test simulates the realistic cloud QPU congestion by injecting random delays of 0.5--2.0\,s into the statevector 
evaluation on odd-numbered, asymmetrical MPI ranks for unpredictable delays in the system, tested through 10 VQE iterations using $H_{2}$.  
Criteria for passing is if all iterations complete without MPI timeout or a deadlock, all returned energies are finite 
with no NaN or infinity from the delays, and lastly if the \texttt{MPI\_Allreduce} synchronization is maintained despite the variable per-rank timing. 
The energy spread across iterations is then recorded for diagnostic purposes. 

\subsubsection{Drop-Out Recovery}
This final test simulates a midrun failure in the stack by deliberately deleting the most recent checkpoint after 10 iterations, 
forcing a recovery from the earlier checkpoint at iteration 5. The procedure here is:
\begin{enumerate}
    \item Run $H_{2}$ VQE for 10 iterations (checkpoints are saved at iterations 5 and 10).
    \item Delete the iteration-10 checkpoint, simulating data loss from a system crash from corruption or data logistics.
    \item Restart from the iteration 5 checkpoint and runs 10 additional iterations (6--15), ensuring the stack can detect the missing checkpoint and automatically falls back to iteration 5.
    \item Verify the restarted optimization continues from the correct $\theta$ state and SPSA schedule position at $k=6$ not $k=1$.
\end{enumerate}

This test passes if the stack detects the missing iteration 10 checkpoint and falls back to the saved iteration 5, if 
the post recovery energy trajectory continues descending (the optimizer continues improving), and lastly if  after 15 iterations the final energy
is lower than the energy documented at iteration 5, confirming that no optimization progress was permanently lost during these simulated crashes.